% Options for packages loaded elsewhere
% Options for packages loaded elsewhere
\PassOptionsToPackage{unicode}{hyperref}
\PassOptionsToPackage{hyphens}{url}
\PassOptionsToPackage{dvipsnames,svgnames,x11names}{xcolor}
%
\documentclass[
  korean,
]{article}
\usepackage{xcolor}
\usepackage[margin=1in]{geometry}
\usepackage{amsmath,amssymb}
\setcounter{secnumdepth}{5}
\usepackage{iftex}
\ifPDFTeX
  \usepackage[T1]{fontenc}
  \usepackage[utf8]{inputenc}
  \usepackage{textcomp} % provide euro and other symbols
\else % if luatex or xetex
  \usepackage{unicode-math} % this also loads fontspec
  \defaultfontfeatures{Scale=MatchLowercase}
  \defaultfontfeatures[\rmfamily]{Ligatures=TeX,Scale=1}
\fi
\usepackage{lmodern}
\ifPDFTeX\else
  % xetex/luatex font selection
  \setmainfont[]{Noto Sans KR}
  \setmonofont[]{JetBrains Mono}
\fi
% Use upquote if available, for straight quotes in verbatim environments
\IfFileExists{upquote.sty}{\usepackage{upquote}}{}
\IfFileExists{microtype.sty}{% use microtype if available
  \usepackage[]{microtype}
  \UseMicrotypeSet[protrusion]{basicmath} % disable protrusion for tt fonts
}{}
\makeatletter
\@ifundefined{KOMAClassName}{% if non-KOMA class
  \IfFileExists{parskip.sty}{%
    \usepackage{parskip}
  }{% else
    \setlength{\parindent}{0pt}
    \setlength{\parskip}{6pt plus 2pt minus 1pt}}
}{% if KOMA class
  \KOMAoptions{parskip=half}}
\makeatother
% Make \paragraph and \subparagraph free-standing
\makeatletter
\ifx\paragraph\undefined\else
  \let\oldparagraph\paragraph
  \renewcommand{\paragraph}{
    \@ifstar
      \xxxParagraphStar
      \xxxParagraphNoStar
  }
  \newcommand{\xxxParagraphStar}[1]{\oldparagraph*{#1}\mbox{}}
  \newcommand{\xxxParagraphNoStar}[1]{\oldparagraph{#1}\mbox{}}
\fi
\ifx\subparagraph\undefined\else
  \let\oldsubparagraph\subparagraph
  \renewcommand{\subparagraph}{
    \@ifstar
      \xxxSubParagraphStar
      \xxxSubParagraphNoStar
  }
  \newcommand{\xxxSubParagraphStar}[1]{\oldsubparagraph*{#1}\mbox{}}
  \newcommand{\xxxSubParagraphNoStar}[1]{\oldsubparagraph{#1}\mbox{}}
\fi
\makeatother

\usepackage{color}
\usepackage{fancyvrb}
\newcommand{\VerbBar}{|}
\newcommand{\VERB}{\Verb[commandchars=\\\{\}]}
\DefineVerbatimEnvironment{Highlighting}{Verbatim}{commandchars=\\\{\}}
% Add ',fontsize=\small' for more characters per line
\usepackage{framed}
\definecolor{shadecolor}{RGB}{241,243,245}
\newenvironment{Shaded}{\begin{snugshade}}{\end{snugshade}}
\newcommand{\AlertTok}[1]{\textcolor[rgb]{0.68,0.00,0.00}{#1}}
\newcommand{\AnnotationTok}[1]{\textcolor[rgb]{0.37,0.37,0.37}{#1}}
\newcommand{\AttributeTok}[1]{\textcolor[rgb]{0.40,0.45,0.13}{#1}}
\newcommand{\BaseNTok}[1]{\textcolor[rgb]{0.68,0.00,0.00}{#1}}
\newcommand{\BuiltInTok}[1]{\textcolor[rgb]{0.00,0.23,0.31}{#1}}
\newcommand{\CharTok}[1]{\textcolor[rgb]{0.13,0.47,0.30}{#1}}
\newcommand{\CommentTok}[1]{\textcolor[rgb]{0.37,0.37,0.37}{#1}}
\newcommand{\CommentVarTok}[1]{\textcolor[rgb]{0.37,0.37,0.37}{\textit{#1}}}
\newcommand{\ConstantTok}[1]{\textcolor[rgb]{0.56,0.35,0.01}{#1}}
\newcommand{\ControlFlowTok}[1]{\textcolor[rgb]{0.00,0.23,0.31}{\textbf{#1}}}
\newcommand{\DataTypeTok}[1]{\textcolor[rgb]{0.68,0.00,0.00}{#1}}
\newcommand{\DecValTok}[1]{\textcolor[rgb]{0.68,0.00,0.00}{#1}}
\newcommand{\DocumentationTok}[1]{\textcolor[rgb]{0.37,0.37,0.37}{\textit{#1}}}
\newcommand{\ErrorTok}[1]{\textcolor[rgb]{0.68,0.00,0.00}{#1}}
\newcommand{\ExtensionTok}[1]{\textcolor[rgb]{0.00,0.23,0.31}{#1}}
\newcommand{\FloatTok}[1]{\textcolor[rgb]{0.68,0.00,0.00}{#1}}
\newcommand{\FunctionTok}[1]{\textcolor[rgb]{0.28,0.35,0.67}{#1}}
\newcommand{\ImportTok}[1]{\textcolor[rgb]{0.00,0.46,0.62}{#1}}
\newcommand{\InformationTok}[1]{\textcolor[rgb]{0.37,0.37,0.37}{#1}}
\newcommand{\KeywordTok}[1]{\textcolor[rgb]{0.00,0.23,0.31}{\textbf{#1}}}
\newcommand{\NormalTok}[1]{\textcolor[rgb]{0.00,0.23,0.31}{#1}}
\newcommand{\OperatorTok}[1]{\textcolor[rgb]{0.37,0.37,0.37}{#1}}
\newcommand{\OtherTok}[1]{\textcolor[rgb]{0.00,0.23,0.31}{#1}}
\newcommand{\PreprocessorTok}[1]{\textcolor[rgb]{0.68,0.00,0.00}{#1}}
\newcommand{\RegionMarkerTok}[1]{\textcolor[rgb]{0.00,0.23,0.31}{#1}}
\newcommand{\SpecialCharTok}[1]{\textcolor[rgb]{0.37,0.37,0.37}{#1}}
\newcommand{\SpecialStringTok}[1]{\textcolor[rgb]{0.13,0.47,0.30}{#1}}
\newcommand{\StringTok}[1]{\textcolor[rgb]{0.13,0.47,0.30}{#1}}
\newcommand{\VariableTok}[1]{\textcolor[rgb]{0.07,0.07,0.07}{#1}}
\newcommand{\VerbatimStringTok}[1]{\textcolor[rgb]{0.13,0.47,0.30}{#1}}
\newcommand{\WarningTok}[1]{\textcolor[rgb]{0.37,0.37,0.37}{\textit{#1}}}

\usepackage{longtable,booktabs,array}
\newcounter{none} % for unnumbered tables
\usepackage{calc} % for calculating minipage widths
% Correct order of tables after \paragraph or \subparagraph
\usepackage{etoolbox}
\makeatletter
\patchcmd\longtable{\par}{\if@noskipsec\mbox{}\fi\par}{}{}
\makeatother
% Allow footnotes in longtable head/foot
\IfFileExists{footnotehyper.sty}{\usepackage{footnotehyper}}{\usepackage{footnote}}
\makesavenoteenv{longtable}
\usepackage{graphicx}
\makeatletter
\newsavebox\pandoc@box
\newcommand*\pandocbounded[1]{% scales image to fit in text height/width
  \sbox\pandoc@box{#1}%
  \Gscale@div\@tempa{\textheight}{\dimexpr\ht\pandoc@box+\dp\pandoc@box\relax}%
  \Gscale@div\@tempb{\linewidth}{\wd\pandoc@box}%
  \ifdim\@tempb\p@<\@tempa\p@\let\@tempa\@tempb\fi% select the smaller of both
  \ifdim\@tempa\p@<\p@\scalebox{\@tempa}{\usebox\pandoc@box}%
  \else\usebox{\pandoc@box}%
  \fi%
}
% Set default figure placement to htbp
\def\fps@figure{htbp}
\makeatother



\ifLuaTeX
\usepackage[bidi=basic,shorthands=off,provide=*]{babel}
\else
\usepackage[bidi=default,shorthands=off,provide=*]{babel}
\fi
\ifPDFTeX
\else
\babelfont{rm}[]{Noto Sans KR}
\fi
\ifLuaTeX
  \usepackage{selnolig} % disable illegal ligatures
\fi


\setlength{\emergencystretch}{3em} % prevent overfull lines

\providecommand{\tightlist}{%
  \setlength{\itemsep}{0pt}\setlength{\parskip}{0pt}}



 


\makeatletter
\@ifpackageloaded{tcolorbox}{}{\usepackage[skins,breakable]{tcolorbox}}
\@ifpackageloaded{fontawesome5}{}{\usepackage{fontawesome5}}
\definecolor{quarto-callout-color}{HTML}{909090}
\definecolor{quarto-callout-note-color}{HTML}{0758E5}
\definecolor{quarto-callout-important-color}{HTML}{CC1914}
\definecolor{quarto-callout-warning-color}{HTML}{EB9113}
\definecolor{quarto-callout-tip-color}{HTML}{00A047}
\definecolor{quarto-callout-caution-color}{HTML}{FC5300}
\definecolor{quarto-callout-color-frame}{HTML}{acacac}
\definecolor{quarto-callout-note-color-frame}{HTML}{4582ec}
\definecolor{quarto-callout-important-color-frame}{HTML}{d9534f}
\definecolor{quarto-callout-warning-color-frame}{HTML}{f0ad4e}
\definecolor{quarto-callout-tip-color-frame}{HTML}{02b875}
\definecolor{quarto-callout-caution-color-frame}{HTML}{fd7e14}
\makeatother
\makeatletter
\@ifpackageloaded{caption}{}{\usepackage{caption}}
\AtBeginDocument{%
\ifdefined\contentsname
  \renewcommand*\contentsname{목차}
\else
  \newcommand\contentsname{목차}
\fi
\ifdefined\listfigurename
  \renewcommand*\listfigurename{그림 목록}
\else
  \newcommand\listfigurename{그림 목록}
\fi
\ifdefined\listtablename
  \renewcommand*\listtablename{표 목록}
\else
  \newcommand\listtablename{표 목록}
\fi
\ifdefined\figurename
  \renewcommand*\figurename{그림}
\else
  \newcommand\figurename{그림}
\fi
\ifdefined\tablename
  \renewcommand*\tablename{표}
\else
  \newcommand\tablename{표}
\fi
}
\@ifpackageloaded{float}{}{\usepackage{float}}
\floatstyle{ruled}
\@ifundefined{c@chapter}{\newfloat{codelisting}{h}{lop}}{\newfloat{codelisting}{h}{lop}[chapter]}
\floatname{codelisting}{목록}
\newcommand*\listoflistings{\listof{codelisting}{코드 목록}}
\makeatother
\makeatletter
\makeatother
\makeatletter
\@ifpackageloaded{caption}{}{\usepackage{caption}}
\@ifpackageloaded{subcaption}{}{\usepackage{subcaption}}
\makeatother
\usepackage{bookmark}
\IfFileExists{xurl.sty}{\usepackage{xurl}}{} % add URL line breaks if available
\urlstyle{same}
\hypersetup{
  pdftitle={Transformer: 토큰{,} 어텐션{,} 위치 인코딩},
  pdflang={ko},
  colorlinks=true,
  linkcolor={blue},
  filecolor={Maroon},
  citecolor={Blue},
  urlcolor={Blue},
  pdfcreator={LaTeX via pandoc}}


\title{Transformer: 토큰, 어텐션, 위치 인코딩}
\usepackage{etoolbox}
\makeatletter
\providecommand{\subtitle}[1]{% add subtitle to \maketitle
  \apptocmd{\@title}{\par {\large #1 \par}}{}{}
}
\makeatother
\subtitle{Attention as Dynamic Pooling --- 고정 가중치를 넘어 입력에
적응하는 가중치로}
\author{}
\date{2026-04-27}
\begin{document}
\maketitle

\renewcommand*\contentsname{목차}
{
\setcounter{tocdepth}{3}
\tableofcontents
}

\section{한눈에 보기}\label{sec-overview}

\begin{tcolorbox}[enhanced jigsaw, arc=.35mm, bottomrule=.15mm, bottomtitle=1mm, breakable, colback=white, colbacktitle=quarto-callout-tip-color!10!white, colframe=quarto-callout-tip-color-frame, coltitle=black, left=2mm, leftrule=.75mm, opacityback=0, opacitybacktitle=0.6, rightrule=.15mm, title=\textcolor{quarto-callout-tip-color}{\faLightbulb}\hspace{0.5em}{핵심 한 줄}, titlerule=0mm, toprule=.15mm, toptitle=1mm]

\textbf{Transformer는 ``고정 가중치(fixed weight)'' 대신 ``입력에 따라
동적으로 바뀌는 가중치(input-dependent weight)''를 쓰는
시퀀스-투-시퀀스(seq-to-seq) 모델이다.} 그 가중치는 데이터를 직접
가중하지 않고, \textbf{데이터끼리의 유사도(similarity)에 부여}된다.

\end{tcolorbox}

Transformer를 떠받치는 세 가지 기둥은 다음과 같다.

{\def\LTcaptype{none} % do not increment counter
\begin{longtable}[]{@{}
  >{\raggedright\arraybackslash}p{(\linewidth - 4\tabcolsep) * \real{0.3333}}
  >{\raggedright\arraybackslash}p{(\linewidth - 4\tabcolsep) * \real{0.3333}}
  >{\raggedright\arraybackslash}p{(\linewidth - 4\tabcolsep) * \real{0.3333}}@{}}
\toprule\noalign{}
\begin{minipage}[b]{\linewidth}\raggedright
구성요소
\end{minipage} & \begin{minipage}[b]{\linewidth}\raggedright
한 줄 정의
\end{minipage} & \begin{minipage}[b]{\linewidth}\raggedright
직관
\end{minipage} \\
\midrule\noalign{}
\endhead
\bottomrule\noalign{}
\endlastfoot
\textbf{토큰 (Token)} & 정보를 캡슐화한 벡터(뉴런 묶음) & 처리의 최소
단위를 ``뉴런''에서 ``벡터 묶음''으로 격상 \\
\textbf{어텐션 (Attention)} & 유사도에 학습된 가중치를 부여하는 동적
풀링 & 입력을 보고 가중치를 바꿔서 다시 입력에 적용 \\
\textbf{위치 인코딩 (Positional Encoding)} & 토큰에 위치 주소를 더하는
코드 & 순서를 모르는 모델에게 ``몇 번째 토큰인지''를 알려줌 \\
\end{longtable}
}

이 셋 중 \textbf{어텐션이 가장 중요}하다. 나머지 둘은 어텐션을 실제
데이터에 적용 가능하게 만드는 보조 장치에 가깝다.

\begin{center}\rule{0.5\linewidth}{0.5pt}\end{center}

\section{배경: 왜 CNN/RNN으로는 부족한가}\label{sec-motivation}

\subsection{시퀀스 모델링의 두 가지
난제}\label{uxc2dcuxd000uxc2a4-uxbaa8uxb378uxb9c1uxc758-uxb450-uxac00uxc9c0-uxb09cuxc81c}

지금까지 다룬 데이터는 길이가 고정된 ``스냅샷''(이미지, 그래프)이었다.
그러나 언어, 시계열(주가), 비디오, 오디오는 길이가 가변적인
\textbf{시퀀스(sequence)} 다. 시퀀스 인코딩에는 두 가지 핵심 난제가
있다.

\begin{enumerate}
\def\labelenumi{\arabic{enumi}.}
\tightlist
\item
  \textbf{장기 의존성 (long-term dependency)}: 멀리 떨어진 정보를 오래
  기억하기 어렵다.
\item
  \textbf{문맥 (context)}: 같은 토큰이라도 주변 문맥에 따라 의미가
  달라진다. (예: ``Buffalo buffalo Buffalo buffalo \ldots{}'' 문장처럼
  동일 단어가 문맥상 명사/동사/지명으로 갈린다.)
\end{enumerate}

RNN은 은닉 상태(hidden state)를 시작부터 끝까지 전파하여 이 두 난제를
다룬다. Transformer는 이에 대한 \textbf{대안}이며, 현재는 사실상 표준이
되었다.

\subsection{CNN의 한계: VQA로 보는 실패
사례}\label{cnnuxc758-uxd55cuxacc4-vqauxb85c-uxbcf4uxb294-uxc2e4uxd328-uxc0acuxb840}

\begin{tcolorbox}[enhanced jigsaw, arc=.35mm, bottomrule=.15mm, bottomtitle=1mm, breakable, colback=white, colbacktitle=quarto-callout-note-color!10!white, colframe=quarto-callout-note-color-frame, coltitle=black, left=2mm, leftrule=.75mm, opacityback=0, opacitybacktitle=0.6, rightrule=.15mm, title=\textcolor{quarto-callout-note-color}{\faInfo}\hspace{0.5em}{VQA(Visual Question Answering) 예시}, titlerule=0mm, toprule=.15mm, toptitle=1mm]

``이 이미지에 새가 몇 마리 있는가?'', ``오른쪽 위 새와 왼쪽 아래 새가
같은 종인가?'', ``하늘 색은 무엇인가?'' 같은 질문을 생각하자.

\end{tcolorbox}

CNN은 두 가지 설계 가정 위에 서 있다.

\begin{itemize}
\tightlist
\item
  \textbf{국소성 (locality)}: 필터는 인접 영역만 본다.
\item
  \textbf{공간 정상성 (spatial stationarity)}: 같은 필터를 이미지 전역에
  동일하게 적용한다(파라미터 공유).
\end{itemize}

이 가정 때문에 CNN은 다음을 잘 못한다.

\begin{itemize}
\tightlist
\item
  \textbf{개수 세기}: max-pooling은 가장 강하게 반응한 위치만 다음
  층으로 넘긴다. 위치마다 다른 ``새''를 세는 작업과 상충한다.
\item
  \textbf{장거리 관계 비교}: 멀리 떨어진 두 객체(``오른쪽 위 vs 왼쪽
  아래'')의 관계 판단은 국소 필터로는 어렵다.
\item
  \textbf{문맥에 따른 선택적 주목}: CNN은 모든 공간 위치를 동등하게
  훑지만, 질문에 따라 ``특정 영역만'' 주목해야 할 때가 있다.
\end{itemize}

locally-connected network(가중치 공유를 없앤 CNN)는 개수 세기에는
유리하나, 파라미터가 폭증하고 CNN의 이점을 거의 다 잃는다. 그래서
\textbf{다른 종류의 인코딩 방식}이 필요하다 --- 그것이 어텐션이다.

\begin{center}\rule{0.5\linewidth}{0.5pt}\end{center}

\section{토큰 (Token)}\label{sec-token}

\subsection{새로운 데이터
구조}\label{uxc0c8uxb85cuxc6b4-uxb370uxc774uxd130-uxad6cuxc870}

\begin{tcolorbox}[enhanced jigsaw, arc=.35mm, bottomrule=.15mm, bottomtitle=1mm, breakable, colback=white, colbacktitle=quarto-callout-tip-color!10!white, colframe=quarto-callout-tip-color-frame, coltitle=black, left=2mm, leftrule=.75mm, opacityback=0, opacitybacktitle=0.6, rightrule=.15mm, title=\textcolor{quarto-callout-tip-color}{\faLightbulb}\hspace{0.5em}{핵심 한 줄}, titlerule=0mm, toprule=.15mm, toptitle=1mm]

\textbf{토큰은 ``뉴런 하나''가 아니라 ``정보를 캡슐화한 벡터''다.}
Transformer는 뉴런이 아니라 토큰 단위로 연산한다.

\end{tcolorbox}

\begin{itemize}
\tightlist
\item
  GNN의 ``노드 속성(node attribute)''과 본질적으로 같은 발상이다 ---
  정보 단위가 벡터다.
\item
  입력을 \textbf{잘게 쪼개고(chop)}, 각 조각을 벡터로
  \textbf{사영(project)} 하면 토큰이 된다.
\end{itemize}

{\def\LTcaptype{none} % do not increment counter
\begin{longtable}[]{@{}ll@{}}
\toprule\noalign{}
입력 유형 & 토큰화 방법 \\
\midrule\noalign{}
\endhead
\bottomrule\noalign{}
\endlastfoot
이미지 & 패치(patch)로 분할 후 각 패치를 벡터로 사영 (ViT) \\
텍스트 & 단어/서브워드를 임베딩 벡터로 변환 \\
오디오 & 시간 구간(프레임)을 벡터로 변환 \\
\end{longtable}
}

\subsection{한 번 토큰화하면 모든 모달리티가 같은
아키텍처로}\label{uxd55c-uxbc88-uxd1a0uxd070uxd654uxd558uxba74-uxbaa8uxb4e0-uxbaa8uxb2ecuxb9acuxd2f0uxac00-uxac19uxc740-uxc544uxd0a4uxd14duxcc98uxb85c}

이것이 Transformer의 큰 실용적 이점이다. \textbf{모달리티에 상관없이
일단 토큰화하면 동일한 아키텍처를 그대로 쓸 수 있다.}

\begin{quote}
CNN은 정보를 ``인코딩하는 한 가지 방식'', MLP는 정보를 ``혼합하는 한
가지 방식''에 가깝다. Transformer는 입력을 토큰 공간(token space)으로
밀어 넣음으로써, 모달리티 통합의 부담을 ``토큰화'' 단계로 옮긴다.
\end{quote}

토큰 간 선형 결합(linear combination)과 토큰 단위 비선형성(token-wise
nonlinearity, 즉 위치별 MLP)을 거치면서 토큰은 갱신된다.

\begin{center}\rule{0.5\linewidth}{0.5pt}\end{center}

\section{어텐션 (Attention)}\label{sec-attention}

여기가 이 글의 심장이다.

\subsection{고정 가중치 vs
어텐션}\label{uxace0uxc815-uxac00uxc911uxce58-vs-uxc5b4uxd150uxc158}

\begin{tcolorbox}[enhanced jigsaw, arc=.35mm, bottomrule=.15mm, bottomtitle=1mm, breakable, colback=white, colbacktitle=quarto-callout-important-color!10!white, colframe=quarto-callout-important-color-frame, coltitle=black, left=2mm, leftrule=.75mm, opacityback=0, opacitybacktitle=0.6, rightrule=.15mm, title=\textcolor{quarto-callout-important-color}{\faExclamation}\hspace{0.5em}{가장 중요한 구분}, titlerule=0mm, toprule=.15mm, toptitle=1mm]

\begin{itemize}
\tightlist
\item
  \textbf{보통의 가중치(ordinary weight)}: 입력이 무엇이든
  \textbf{고정된 값}. 데이터로 학습되지만, 추론 시에는 입력과 무관하게
  동일.
\item
  \textbf{어텐션(attention)}: 가중치가 \textbf{입력의 함수}. 입력을 받아
  가중치 값을 바꾸고, 그 바뀐 가중치를 다시 입력에 적용.
\end{itemize}

\end{tcolorbox}

수식으로 보면, 보통의 층은 출력이 \(y = Wx\) 로 \(W\)가 상수다. 어텐션은
\(y = A(x)\,x\) 처럼 \textbf{가중치 자체가 \(x\)의 함수} \(A(x)\) 다.
``입력을 입력으로 받아 가중치를 만들고, 그 가중치를 입력에 다시
적용한다''는 한 단계가 추가될 뿐이다.

\subsection{동적 풀링(dynamic pooling)으로서의
어텐션}\label{uxb3d9uxc801-uxd480uxb9c1dynamic-poolinguxc73cuxb85cuxc11cuxc758-uxc5b4uxd150uxc158}

어텐션은 \textbf{쿼리(query)를 기준으로 한 동적/적응적 풀링}으로 이해할
수 있다.

\begin{itemize}
\tightlist
\item
  max/average pooling은 위치 간 \textbf{관계를 고려하지 않고} 고정된
  연산(최대/평균)을 한다.
\item
  어텐션은 \textbf{쿼리와 키 사이의 관계(유사도)에 따라} 풀링 가중치를
  동적으로 정한다.
\end{itemize}

즉 ``max도 min도 average도 아닌, 관계 기반 풀링''이다.

\subsection{쿼리·키·값 (Query, Key,
Value)}\label{uxcffcuxb9acuxd0a4uxac12-query-key-value}

세 가지 역할을 정의한다.

\begin{itemize}
\tightlist
\item
  \textbf{Query (\(q\))}: 인코딩(임베딩)을 얻고자 하는 대상. 내가
  ``물어보는'' 것.
\item
  \textbf{Key (\(k\))}: 문맥을 이루는 다른 정보들. 쿼리와의 유사도를 잴
  상대.
\item
  \textbf{Value (\(v\))}: 실제로 가져와서 합칠 정보. 연산의 결과로
  출력에 반영됨.
\end{itemize}

\begin{tcolorbox}[enhanced jigsaw, arc=.35mm, bottomrule=.15mm, bottomtitle=1mm, breakable, colback=white, colbacktitle=quarto-callout-note-color!10!white, colframe=quarto-callout-note-color-frame, coltitle=black, left=2mm, leftrule=.75mm, opacityback=0, opacitybacktitle=0.6, rightrule=.15mm, title=\textcolor{quarto-callout-note-color}{\faInfo}\hspace{0.5em}{언어 예시: ``I like ice cream''}, titlerule=0mm, toprule=.15mm, toptitle=1mm]

``like''의 임베딩을 얻고 싶다면, \textbf{``like''가 쿼리}, 문장의 모든
단어(``I'', ``like'', ``ice'', ``cream'')가 \textbf{키}가 된다. 쿼리와
각 키의 유사도를 재고, 그 유사도로 값들을 가중합한다.

\end{tcolorbox}

\begin{tcolorbox}[enhanced jigsaw, arc=.35mm, bottomrule=.15mm, bottomtitle=1mm, breakable, colback=white, colbacktitle=quarto-callout-note-color!10!white, colframe=quarto-callout-note-color-frame, coltitle=black, left=2mm, leftrule=.75mm, opacityback=0, opacitybacktitle=0.6, rightrule=.15mm, title=\textcolor{quarto-callout-note-color}{\faInfo}\hspace{0.5em}{VQA 예시 (cross-attention)}, titlerule=0mm, toprule=.15mm, toptitle=1mm]

``이 임팔라의 색은?''이라는 \textbf{질문 문장이 쿼리}, \textbf{이미지
패치들이 키/값}이 된다. 쿼리와 비슷한 패치(임팔라 영역)에 큰 가중치가
실리고, 그 영역의 값으로부터 답의 단서를 얻는다. 쿼리와 키가 서로 다른
출처일 때 이를 \textbf{cross-attention}이라 한다.

\end{tcolorbox}

문제 설정(번역, VQA, 비디오 생성 등)에 따라 ``무엇을 쿼리/키/값으로
둘지''만 정의하면, 다양한 문제를 하나의 형식으로 환원할 수 있다.

\subsection{유사도는 왜 코사인 유사도인가}\label{sec-cosine}

쿼리 \(q\)와 키 \(k\)의 유사도를 \textbf{내적(dot product)} 으로 잰다.
그 의미를 고등학교 벡터로 풀어보자.

\[
q \cdot k = \|q\|\,\|k\|\cos\theta
\]

\begin{itemize}
\tightlist
\item
  \(\|q\|,\|k\|\) 는 크기(magnitude), \(\theta\) 는 두 벡터 사이의 각.
\item
  두 벡터를 단위 벡터로 정규화하면(또는 크기를 무시하면) 내적은
  \(\cos\theta\) 그 자체다.
\end{itemize}

\[
\text{cosine similarity} = \frac{q\cdot k}{\|q\|\,\|k\|} = \cos\theta
\]

기하학적으로 \(\cos\theta\) 는 한 벡터를 다른 벡터 방향으로
\textbf{사영(projection)} 한 길이다. 즉:

\begin{itemize}
\tightlist
\item
  \(\theta = 90°\) (직교) → \(\cos\theta = 0\): 무관.
\item
  \(\theta = 0°\) (같은 방향) → \(\cos\theta = 1\): 최대 유사.
\end{itemize}

\begin{tcolorbox}[enhanced jigsaw, arc=.35mm, bottomrule=.15mm, bottomtitle=1mm, breakable, colback=white, colbacktitle=quarto-callout-tip-color!10!white, colframe=quarto-callout-tip-color-frame, coltitle=black, left=2mm, leftrule=.75mm, opacityback=0, opacitybacktitle=0.6, rightrule=.15mm, title=\textcolor{quarto-callout-tip-color}{\faLightbulb}\hspace{0.5em}{직관}, titlerule=0mm, toprule=.15mm, toptitle=1mm]

코사인 유사도는 \textbf{크기를 무시하고 ``방향이 얼마나 같은 쪽을
향하는가''} 만 본다. 이를 \textbf{각거리(angular distance)} 라고도
부른다. 어텐션이 묻는 것은 ``쿼리와 키가 같은 방향을 향하느냐''이다.

\end{tcolorbox}

\subsection{소프트맥스: 차이를
증폭하기}\label{uxc18cuxd504uxd2b8uxb9e5uxc2a4-uxcc28uxc774uxb97c-uxc99duxd3eduxd558uxae30}

유사도 점수 벡터 \(s = [s_1, \dots, s_n]\) 에
\textbf{소프트맥스(softmax)} 를 적용한다.

\[
a_j = \frac{\exp(s_j)}{\sum_{l} \exp(s_l)}
\]

소프트맥스의 역할은 \textbf{최댓값과 나머지의 차이를 증폭}하여, 합이 1인
가중치 분포로 만드는 것이다(부드러운 argmax). 이렇게 얻은 \(a\) 가 바로
\textbf{어텐션 가중치}다.

\subsection{\texorpdfstring{값 변환 행렬 \(W_v\) 와 학습
대상}{값 변환 행렬 W\_v 와 학습 대상}}\label{uxac12-uxbcc0uxd658-uxd589uxb82c-w_v-uxc640-uxd559uxc2b5-uxb300uxc0c1}

지금까지 \(q, k\) 는 유사도(=어텐션 강도)를 정하는 데 쓰였다. 그렇다면
실제로 합쳐질 \textbf{값(value)} 은? 입력 토큰 \(x\) 를 그대로 쓰지
않고, 학습 가능한 행렬 \(W_v\) 로 한 번 더 변환한다. 이는 주어진 과제에
맞게 값 공간을 회전/스케일할 자유도를 준다.

정리하면, 어텐션이 학습하는 것은 \textbf{세 개의 사영 행렬}
\(W_q, W_k, W_v\) 다.

\[
Q = X W_q,\quad K = X W_k,\quad V = X W_v
\]

\begin{tcolorbox}[enhanced jigsaw, arc=.35mm, bottomrule=.15mm, bottomtitle=1mm, breakable, colback=white, colbacktitle=quarto-callout-important-color!10!white, colframe=quarto-callout-important-color-frame, coltitle=black, left=2mm, leftrule=.75mm, opacityback=0, opacitybacktitle=0.6, rightrule=.15mm, title=\textcolor{quarto-callout-important-color}{\faExclamation}\hspace{0.5em}{어텐션의 본질}, titlerule=0mm, toprule=.15mm, toptitle=1mm]

어텐션은 ``\textbf{정보 자체에 고정 가중치를 거는 것}''이 아니라,
``\textbf{정보들 사이의 유사도에 학습된 가중치를 거는 것}''이다. 어떤
유사도가 중요한지는 데이터와 과제에 따라 달라지며, 그 판단을
\(W_q, W_k, W_v\) 가 학습한다.

\end{tcolorbox}

\subsection{스케일드 닷-프로덕트 어텐션 (보충)}\label{sec-scaled}

전체 연산을 행렬로 한 번에 쓰면 다음과 같다. (이 정규화 계수
\(\sqrt{d_k}\) 는 위 직관 설명에서는 생략되지만, 실제 구현에서 반드시
등장하므로 보충한다 --- d2l.ai의 \emph{Attention Scoring Functions} 절
참고.)

\[
\boxed{\ \text{Attention}(Q, K, V) = \mathrm{softmax}\!\left(\frac{QK^\top}{\sqrt{d_k}}\right) V\ }
\]

\begin{tcolorbox}[enhanced jigsaw, arc=.35mm, bottomrule=.15mm, bottomtitle=1mm, breakable, colback=white, colbacktitle=quarto-callout-note-color!10!white, colframe=quarto-callout-note-color-frame, coltitle=black, left=2mm, leftrule=.75mm, opacityback=0, opacitybacktitle=0.6, rightrule=.15mm, title=\textcolor{quarto-callout-note-color}{\faInfo}\hspace{0.5em}{왜 \(\sqrt{d_k}\) 로 나누는가}, titlerule=0mm, toprule=.15mm, toptitle=1mm]

\(q, k\) 의 각 성분이 평균 0, 분산 1로 독립이라면, 내적
\(q\cdot k = \sum_{i=1}^{d_k} q_i k_i\) 의 분산은 \(d_k\) 가 된다.
\(d_k\) 가 크면 점수의 크기가 커져 소프트맥스가 한쪽으로
과포화(saturation)되고, 그 영역에서 기울기가 거의 0이 된다. 표준편차
\(\sqrt{d_k}\) 로 나누면 점수 분산을 1 수준으로 되돌려 학습을
안정화한다.

\end{tcolorbox}

\begin{center}\rule{0.5\linewidth}{0.5pt}\end{center}

\section{셀프 어텐션 (Self-Attention)}\label{sec-self}

\begin{tcolorbox}[enhanced jigsaw, arc=.35mm, bottomrule=.15mm, bottomtitle=1mm, breakable, colback=white, colbacktitle=quarto-callout-tip-color!10!white, colframe=quarto-callout-tip-color-frame, coltitle=black, left=2mm, leftrule=.75mm, opacityback=0, opacitybacktitle=0.6, rightrule=.15mm, title=\textcolor{quarto-callout-tip-color}{\faLightbulb}\hspace{0.5em}{핵심 한 줄}, titlerule=0mm, toprule=.15mm, toptitle=1mm]

\textbf{쿼리·키·값이 모두 같은 출처에서 나오면 셀프 어텐션}이다. 각
토큰이 ``자기 자신을 쿼리로, 주변 모두(자신 포함)를 키/값으로'' 삼아,
문맥을 반영한 표현으로 자신을 정련(refine)한다.

\end{tcolorbox}

\begin{itemize}
\tightlist
\item
  예: 한 이미지 패치의 임베딩을 얻을 때, \textbf{그 패치를 쿼리},
  \textbf{같은 이미지의 모든 패치를 키/값}으로 둔다.
\item
  결과: 각 토큰은 이웃 토큰 정보를 반영한, \textbf{의미적으로 더
  풍부한(semantic)} 임베딩을 갖게 된다.
\end{itemize}

\subsection{행렬 연산}\label{uxd589uxb82c-uxc5f0uxc0b0}

\(X \in \mathbb{R}^{n\times d}\) ( \(n\) 토큰, 차원 \(d\) )에 대해:

\[
Q=XW_q,\ K=XW_k,\ V=XW_v;\qquad
\text{Out} = \mathrm{softmax}\!\left(\tfrac{QK^\top}{\sqrt{d_k}}\right)V
\]

연산 순서는 cross-attention과 동일하며, \(Q,K,V\) 가 같은 \(X\) 에서
나온다는 점만 다르다.

\subsection{DINO: 패치 값 정련의
시각화}\label{dino-uxd328uxce58-uxac12-uxc815uxb828uxc758-uxc2dcuxac01uxd654}

자기지도학습(self-supervised learning) 모델 \textbf{DINO}(Caron et al.,
ICCV 2021)에서, 한 작은 패치를 쿼리로 두고 어텐션 점수를 시각화하면
\textbf{의미적으로 같은 객체(개) 영역이 밝게}, 무관한 배경은 어둡게
나타난다. 셀프 어텐션이 의미 단위로 패치를 묶고 있음을 보여준다.

\begin{center}\rule{0.5\linewidth}{0.5pt}\end{center}

\section{마스크드 어텐션 (Masked / Causal Attention)}\label{sec-mask}

\begin{tcolorbox}[enhanced jigsaw, arc=.35mm, bottomrule=.15mm, bottomtitle=1mm, breakable, colback=white, colbacktitle=quarto-callout-tip-color!10!white, colframe=quarto-callout-tip-color-frame, coltitle=black, left=2mm, leftrule=.75mm, opacityback=0, opacitybacktitle=0.6, rightrule=.15mm, title=\textcolor{quarto-callout-tip-color}{\faLightbulb}\hspace{0.5em}{핵심 한 줄}, titlerule=0mm, toprule=.15mm, toptitle=1mm]

\textbf{어텐션 행렬의 일부를 0으로 막아, 특정 토큰이 특정 토큰에
주목하지 못하게 한다.} 디코더에서는 ``미래 토큰''을 막는 데 쓴다(causal
attention).

\end{tcolorbox}

\subsection{왜 미래를
막는가}\label{uxc65c-uxbbf8uxb798uxb97c-uxb9c9uxb294uxac00}

기계 번역(seq-to-seq)을 생각하자. 인코더는 입력 문장 전체를 보지만,
\textbf{디코더는 단어를 하나씩 생성}한다. \(t\) 시점에서 \(t+1\) 이후의
단어는 \textbf{아직 생성되지 않았다.}

\begin{itemize}
\tightlist
\item
  만약 미래 토큰에 어텐션을 허용하면, 학습 시 디코더가 정답(미래 단어)을
  ``엿보게'' 되고, 추론 시에는 존재하지 않는 토큰을 참조하게 된다.
\item
  그 결과 미래 위치의 어텐션 점수는 무의미한 값이 되고, 역전파 시
  \textbf{쓸모없는(garbage) 기울기}가 \(W_q, W_k\) 까지 전파되어 학습을
  망친다.
\end{itemize}

해결책: 미래 위치의 점수를 \(-\infty\) (실제로는 매우 큰 음수)로 만들어
소프트맥스 후 가중치가 0이 되게 한다.

\[
s_{ij} \leftarrow
\begin{cases}
s_{ij} & j \le i \quad(\text{과거·현재})\\[4pt]
-\infty & j > i \quad(\text{미래, 마스킹})
\end{cases}
\]

\subsection{다른 용도: 선택적 멀티모달
인코딩}\label{uxb2e4uxb978-uxc6a9uxb3c4-uxc120uxd0dduxc801-uxba40uxd2f0uxbaa8uxb2ec-uxc778uxcf54uxb529}

마스킹은 인과성 외에도, \textbf{어떤 토큰을 문맥에서 배제할지 선택}하는
일반적 장치로 쓸 수 있다. 예컨대 텍스트 스토리에 부합하는 일관된 이미지
시퀀스를 생성할 때, 어텐션 마스크를 설계해 관련 토큰만 선택적으로
인코딩할 수 있다(Ahn et al., \emph{Story visualization by online text
augmentation with context memory}, ICCV 2023).

\begin{center}\rule{0.5\linewidth}{0.5pt}\end{center}

\section{멀티헤드 셀프 어텐션 (Multi-Head
Self-Attention)}\label{sec-mha}

\begin{tcolorbox}[enhanced jigsaw, arc=.35mm, bottomrule=.15mm, bottomtitle=1mm, breakable, colback=white, colbacktitle=quarto-callout-tip-color!10!white, colframe=quarto-callout-tip-color-frame, coltitle=black, left=2mm, leftrule=.75mm, opacityback=0, opacitybacktitle=0.6, rightrule=.15mm, title=\textcolor{quarto-callout-tip-color}{\faLightbulb}\hspace{0.5em}{핵심 한 줄}, titlerule=0mm, toprule=.15mm, toptitle=1mm]

\textbf{하나의 \((W_q, W_k, W_v)\) 쌍은 관계의 한 측면(aspect)만
포착한다. 여러 측면을 보려면 여러 쌍을 병렬로 둔다.} 이것이 CNN에서 여러
필터를 학습하는 것과 같은 동기다.

\end{tcolorbox}

\begin{itemize}
\tightlist
\item
  CNN: 한 이미지에 대해 여러 필터를 학습 → 각기 다른 공간 패턴 포착.
\item
  Transformer: 여러 헤드
  \(\{(W_q^{(h)}, W_k^{(h)}, W_v^{(h)})\}_{h=1}^{H}\) 를 병렬로 → 각기
  다른 관계 측면 포착.
\end{itemize}

각 헤드의 출력을 이어 붙인(concatenate) 뒤 출력 사영 \(W_o\) 로 섞는다.

\[
\text{head}_h = \text{Attention}(XW_q^{(h)}, XW_k^{(h)}, XW_v^{(h)})
\] \[
\text{MHA}(X) = \big[\text{head}_1; \dots; \text{head}_H\big]\,W_o
\]

\begin{tcolorbox}[enhanced jigsaw, arc=.35mm, bottomrule=.15mm, bottomtitle=1mm, breakable, colback=white, colbacktitle=quarto-callout-warning-color!10!white, colframe=quarto-callout-warning-color-frame, coltitle=black, left=2mm, leftrule=.75mm, opacityback=0, opacitybacktitle=0.6, rightrule=.15mm, title=\textcolor{quarto-callout-warning-color}{\faExclamationTriangle}\hspace{0.5em}{솔직한 한계}, titlerule=0mm, toprule=.15mm, toptitle=1mm]

서로 다른 헤드 사이에 ``지식 중복 방지''나 ``정보 소멸 방지''를 위한
\textbf{명시적 메커니즘은 없다.} 헤드들이 알아서 다른 측면을 학습하리라
\textbf{암묵적으로 가정}할 뿐이다. 헤드 간 상호작용은 다음 층의 혼합
단계에서 암묵적으로만 일어난다. 어텐션을 빼면 성능이 의미 있게(그러나
치명적이지는 않게) 떨어진다 --- 우리는 많은 추측(conjecture)은 갖고
있으나 정확한 이유는 모른다.

\end{tcolorbox}

여러 헤드(헤드 내 병렬)에 더해, 블록을 \textbf{여러 층으로
쌓는다(cascade)}. CNN에서 합성곱 층을 쌓는 것과 같다.

\begin{center}\rule{0.5\linewidth}{0.5pt}\end{center}

\section{MLP vs Transformer, 그리고 순열 등변성}\label{sec-equivariance}

MLP와 Transformer의 결정적 차이: \textbf{Transformer의 어텐션은 입력의
함수}라는 점이다(\(A(x)\)). MLP의 가중치는 입력과 무관하게 고정이다.

\subsection{등변성(equivariance)과
불변성(invariance)}\label{uxb4f1uxbcc0uxc131equivarianceuxacfc-uxbd88uxbcc0uxc131invariance}

{\def\LTcaptype{none} % do not increment counter
\begin{longtable}[]{@{}
  >{\raggedright\arraybackslash}p{(\linewidth - 4\tabcolsep) * \real{0.3333}}
  >{\raggedright\arraybackslash}p{(\linewidth - 4\tabcolsep) * \real{0.3333}}
  >{\raggedright\arraybackslash}p{(\linewidth - 4\tabcolsep) * \real{0.3333}}@{}}
\toprule\noalign{}
\begin{minipage}[b]{\linewidth}\raggedright
개념
\end{minipage} & \begin{minipage}[b]{\linewidth}\raggedright
정의
\end{minipage} & \begin{minipage}[b]{\linewidth}\raggedright
한 줄 직관
\end{minipage} \\
\midrule\noalign{}
\endhead
\bottomrule\noalign{}
\endlastfoot
\textbf{불변성 (invariance)} & 입력이 바뀌어도 출력이 그대로 & 변화에
\textbf{둔감} \\
\textbf{등변성 (equivariance)} & 입력이 바뀐 만큼 출력도 같이 바뀜 &
변화를 \textbf{반영} \\
\end{longtable}
}

예: 입력을 회전시켰을 때 출력도 같이 회전하면 회전 등변, 출력이 그대로면
회전 불변.

\subsection{순열 등변성과 위치 인코딩의
필요성}\label{uxc21cuxc5f4-uxb4f1uxbcc0uxc131uxacfc-uxc704uxce58-uxc778uxcf54uxb529uxc758-uxd544uxc694uxc131}

\begin{tcolorbox}[enhanced jigsaw, arc=.35mm, bottomrule=.15mm, bottomtitle=1mm, breakable, colback=white, colbacktitle=quarto-callout-important-color!10!white, colframe=quarto-callout-important-color-frame, coltitle=black, left=2mm, leftrule=.75mm, opacityback=0, opacitybacktitle=0.6, rightrule=.15mm, title=\textcolor{quarto-callout-important-color}{\faExclamation}\hspace{0.5em}{문제}, titlerule=0mm, toprule=.15mm, toptitle=1mm]

어텐션(과 MLP)은 \textbf{순열 등변(permutation equivariant)} 이다. 입력
토큰의 순서를 섞으면 출력도 똑같이 섞일 뿐, \textbf{모델은 토큰의 원래
위치를 모른다.}

\end{tcolorbox}

언어와 이미지에서 위치/순서는 의미를 좌우한다(``개가 사람을 물다'' ≠
``사람이 개를 물다''). 따라서 토큰에 \textbf{위치 정보를 명시적으로
주입}해야 한다. 그것이 다음 절의 위치 인코딩이다.

\begin{center}\rule{0.5\linewidth}{0.5pt}\end{center}

\section{위치 인코딩 (Positional Encoding)}\label{sec-pe}

\begin{tcolorbox}[enhanced jigsaw, arc=.35mm, bottomrule=.15mm, bottomtitle=1mm, breakable, colback=white, colbacktitle=quarto-callout-tip-color!10!white, colframe=quarto-callout-tip-color-frame, coltitle=black, left=2mm, leftrule=.75mm, opacityback=0, opacitybacktitle=0.6, rightrule=.15mm, title=\textcolor{quarto-callout-tip-color}{\faLightbulb}\hspace{0.5em}{핵심 한 줄}, titlerule=0mm, toprule=.15mm, toptitle=1mm]

\textbf{각 위치마다 다른 ``주소 벡터''를 입력 토큰에 더한다.} 단, 이
주소는 학습이 잘 되도록 \textbf{조밀(dense)하고 비슷한 크기}여야 한다.

\end{tcolorbox}

\subsection{왜 단순 정수 번호는 안
되는가}\label{uxc65c-uxb2e8uxc21c-uxc815uxc218-uxbc88uxd638uxb294-uxc548-uxb418uxb294uxac00}

가장 단순한 발상은 위치를 정수 \(0,1,2,\dots\) 로 매기는 것이다. 그러나
이진 표현처럼 0이 많은 \textbf{희소(sparse)} 코드는 그 위에 가중치를
학습하기 어렵다(0에 곱해지는 가중치는 사실상 학습되지 않음). 우리는
\textbf{모든 차원이 의미 있는 값을 갖는 조밀한 표현}, 그리고 위치마다
다르되 \textbf{크기가 서로 비슷한} 벡터가 필요하다.

\subsection{사인파 위치 인코딩
(보충)}\label{uxc0acuxc778uxd30c-uxc704uxce58-uxc778uxcf54uxb529-uxbcf4uxcda9}

가장 널리 쓰이는 구현은 서로 다른 주파수의 사인/코사인을 쓰는
것이다(Vaswani et al., 2017; d2l.ai \emph{Self-Attention and Positional
Encoding} 절).

\[
PE_{(pos,\,2i)} = \sin\!\left(\frac{pos}{10000^{\,2i/d}}\right),\qquad
PE_{(pos,\,2i+1)} = \cos\!\left(\frac{pos}{10000^{\,2i/d}}\right)
\]

\begin{itemize}
\tightlist
\item
  모든 성분의 크기가 \([-1, 1]\) 로 균일 → 가중치 학습에 유리.
\item
  위치마다 고유한 패턴(주소) → 위치 구분 가능.
\item
  서로 다른 주파수의 중첩이라 상대 위치 정보도 선형 변환으로 표현 가능.
\end{itemize}

\begin{tcolorbox}[enhanced jigsaw, arc=.35mm, bottomrule=.15mm, bottomtitle=1mm, breakable, colback=white, colbacktitle=quarto-callout-note-color!10!white, colframe=quarto-callout-note-color-frame, coltitle=black, left=2mm, leftrule=.75mm, opacityback=0, opacitybacktitle=0.6, rightrule=.15mm, title=\textcolor{quarto-callout-note-color}{\faInfo}\hspace{0.5em}{비유: 군부대 번호 매기기}, titlerule=0mm, toprule=.15mm, toptitle=1mm]

값 자체는 ``비상식적''으로 보여도(예: 6575번 부대) 학습 가중치
입장에서는 좋은 성질을 갖도록 설계된 코드다. 사인파 PE는 ``인간 눈엔
무의미해 보이지만 기계 학습엔 유용한'' 위치 주소다.

\end{tcolorbox}

위치 인코딩까지 더하면, 멀티헤드 (마스크드) 셀프 어텐션 + 피드포워드가
문맥을 잘 포착하여 적절한 출력 확률을 만든다.

\begin{center}\rule{0.5\linewidth}{0.5pt}\end{center}

\section{전체 아키텍처 정리}\label{sec-arch}

이제 블록을 조립하자. 인코더-디코더 구조는 다음과 같다.

\includegraphics[width=7.47in,height=12.52in]{DL-lecture-260427-notes_files/figure-latex/mermaid-figure-1.png}

핵심 흐름:

\begin{enumerate}
\def\labelenumi{\arabic{enumi}.}
\tightlist
\item
  \textbf{인코더}: (마스크 없는) 셀프 어텐션으로 입력 전체의 문맥을
  인코딩.
\item
  \textbf{디코더}: 먼저 \textbf{마스크드} 셀프 어텐션으로 지금까지
  생성된 토큰만 보고, 이어서 \textbf{크로스 어텐션}으로 인코더의
  문맥(Key/Value)과 섞는다.
\item
  잔차 연결(Add)과 정규화(Norm), 위치별 피드포워드가 각 단계를
  안정화·변환.
\item
  최종 출력은 다음 토큰의 \textbf{확률 분포}로, 이 확률이 다음에 어떤
  토큰을 생성할지의 단서가 된다.
\end{enumerate}

\begin{tcolorbox}[enhanced jigsaw, arc=.35mm, bottomrule=.15mm, bottomtitle=1mm, breakable, colback=white, colbacktitle=quarto-callout-note-color!10!white, colframe=quarto-callout-note-color-frame, coltitle=black, left=2mm, leftrule=.75mm, opacityback=0, opacitybacktitle=0.6, rightrule=.15mm, title=\textcolor{quarto-callout-note-color}{\faInfo}\hspace{0.5em}{Add \& Norm --- 왜 잔차 연결인가}, titlerule=0mm, toprule=.15mm, toptitle=1mm]

\(\text{LayerNorm}(x + \text{Sublayer}(x))\) 형태의 잔차 연결은 ResNet과
같은 이유로 쓰인다. 항등 경로(+I) 덕분에 기울기가 깊은 층까지 보존되어
학습 가능성이 유지된다.

\end{tcolorbox}

\begin{center}\rule{0.5\linewidth}{0.5pt}\end{center}

\section{응용 (Application)}\label{sec-application}

\subsection{언어 모델 (Language
Model)}\label{uxc5b8uxc5b4-uxbaa8uxb378-language-model}

언어 모델은 \textbf{이전 단어들이 주어졌을 때 다음 단어의 확률}을
예측하는 NLP 과제다. 흔한 가정은 ``다음 단어 확률이 고정 크기
윈도(window) 내 이전 단어들에만 의존한다''는 것이다. 규모가 매우 크면
\textbf{대규모 언어 모델(LLM)} 이라 부른다(GPT, Claude, HyperClova X
등).

언어 모델링은 사실상 NLP의 \textbf{메타 과제}다. 좋은 언어 모델을
학습하면 감성 분석, 개체명 인식(NER), 공참조 해소, 질의응답 등 수많은
다운스트림 과제를 끌어낼 수 있다. 그래서 사람들은 모델을 점점 키웠다.

\subsection{자기회귀 모델과
GPT}\label{uxc790uxae30uxd68cuxadc0-uxbaa8uxb378uxacfc-gpt}

\textbf{GPT = Generative Pretrained Transformer.} 이름에
``Transformer''가 들어 있다(GPT가 General Purpose Text가 아니다).

\begin{itemize}
\tightlist
\item
  \textbf{자기회귀(autoregressive)}: 이전까지 생성된 토큰을 입력으로
  다음 토큰을 생성, 이를 반복.
\item
  \textbf{학습}: 인코더-디코더(또는 디코더-온리) 모델을 통째로, 한
  단어씩 뱉도록 학습. \textbf{샘플링(추론)} 도 한 토큰씩 생성.
\item
  \textbf{인과 어텐션(causal attention)}: 미래 토큰을 보지 않으므로,
  자기회귀 생성과 정확히 맞물린다.
\end{itemize}

\begin{tcolorbox}[enhanced jigsaw, arc=.35mm, bottomrule=.15mm, bottomtitle=1mm, breakable, colback=white, colbacktitle=quarto-callout-warning-color!10!white, colframe=quarto-callout-warning-color-frame, coltitle=black, left=2mm, leftrule=.75mm, opacityback=0, opacitybacktitle=0.6, rightrule=.15mm, title=\textcolor{quarto-callout-warning-color}{\faExclamationTriangle}\hspace{0.5em}{자기회귀의 약점}, titlerule=0mm, toprule=.15mm, toptitle=1mm]

\textbf{오류 전파(error propagation)}: 초반에 한 번 틀리면 그 오류가
이후 전체 시퀀스로 번진다. 대안으로 단어 덩어리(chunk)를 한꺼번에
생성하는 방식(예: 일부 모델은 CNN/Diffusion으로 청크 생성)도 연구된다.
다만 초기 LLM(GPT-2/3 등)은 Transformer를 그대로 자기회귀로 썼다.

\end{tcolorbox}

\begin{tcolorbox}[enhanced jigsaw, arc=.35mm, bottomrule=.15mm, bottomtitle=1mm, breakable, colback=white, colbacktitle=quarto-callout-note-color!10!white, colframe=quarto-callout-note-color-frame, coltitle=black, left=2mm, leftrule=.75mm, opacityback=0, opacitybacktitle=0.6, rightrule=.15mm, title=\textcolor{quarto-callout-note-color}{\faInfo}\hspace{0.5em}{직접 만져보기: minGPT}, titlerule=0mm, toprule=.15mm, toptitle=1mm]

A. Karpathy의
\textbf{minGPT}(\url{https://github.com/karpathy/minGPT})는
마스크드(causal) 어텐션이 모두 구현된 최소 GPT다. 소스가 매우 단순해
학습용으로 좋다.

\end{tcolorbox}

\begin{center}\rule{0.5\linewidth}{0.5pt}\end{center}

\section{한계와 효율화 (Limitations)}\label{sec-limitation}

Transformer는 거의 10년 동안(2017년 제안 이후) 표준이었지만 한계가 있다.

\subsection{1) 문맥 윈도 밖은 못
본다}\label{uxbb38uxb9e5-uxc708uxb3c4-uxbc16uxc740-uxbabb-uxbcf8uxb2e4}

키 윈도(예: 12.8만 토큰) 안에서만 \(W_q, W_k, W_v\) 를 계산한다. 윈도 밖
정보는 RNN처럼 전파되지 않으므로 잃어버린다 → 멀리 있는 정보를 물으면
\textbf{환각(hallucination)} 이 늘 수 있다.

\begin{itemize}
\tightlist
\item
  \textbf{부분적 완화: KV 캐시(KV cache)}. 중요한 토큰의 Key/Value를
  저장해 두고 재사용하여, 긴 문맥을 사실상 유지한다. 다만 이는
  \textbf{임시 해법}이며 근본 해결은 아니다.
\item
  KV 캐시를 들고 있어야 하므로 \textbf{메모리 수요가 크다}(GPU
  메모리/대역폭 병목의 한 원인).
\end{itemize}

\subsection{\texorpdfstring{2) 계산 복잡도
\(O(n^2)\)}{2) 계산 복잡도 O(n\^{}2)}}\label{uxacc4uxc0b0-uxbcf5uxc7a1uxb3c4-on2}

윈도 크기 \(n\) 에 대해 어텐션은 \(O(n^2)\) 다. \(QK^\top\) 에서 모든
쿼리-키 쌍의 내적을 한 번, 어텐션 가중치를 \(V\) 에 적용할 때 다시 한 번
--- 윈도 크기에 제곱으로 비례하는 연산이 두 번 들어간다. 윈도가 커질수록
비용이 급증한다.

\subsection{효율적 대안}\label{uxd6a8uxc728uxc801-uxb300uxc548}

\begin{itemize}
\tightlist
\item
  \textbf{효율적 Transformer 계열}(Tay et al., \emph{Efficient
  Transformers: A Survey}, 2022): 희소/저랭크 근사 등으로 \(O(n^2)\)
  완화.
\item
  \textbf{Mamba} 등 상태공간모델(SSM) 계열: 긴 문맥을 더 낮은 복잡도로
  인코딩하려는 시도.
\end{itemize}

\begin{tcolorbox}[enhanced jigsaw, arc=.35mm, bottomrule=.15mm, bottomtitle=1mm, breakable, colback=white, colbacktitle=quarto-callout-tip-color!10!white, colframe=quarto-callout-tip-color-frame, coltitle=black, left=2mm, leftrule=.75mm, opacityback=0, opacitybacktitle=0.6, rightrule=.15mm, title=\textcolor{quarto-callout-tip-color}{\faLightbulb}\hspace{0.5em}{열린 문제}, titlerule=0mm, toprule=.15mm, toptitle=1mm]

현재의 Transformer 아키텍처가 인간이 상상할 수 있는 \textbf{최적 구조는
아니다.} 긴 문맥을 더 효율적으로 인코딩하는 결정론적 함수(아키텍처)를
찾는 일은 여전히 열려 있는 큰 연구 주제다.

\end{tcolorbox}

\begin{center}\rule{0.5\linewidth}{0.5pt}\end{center}

\section{무선통신 관점 비유}\label{sec-wireless}

신호처리·무선통신 배경에서 어텐션은 익숙한 개념의 재조합으로 읽힌다.

{\def\LTcaptype{none} % do not increment counter
\begin{longtable}[]{@{}
  >{\raggedright\arraybackslash}p{(\linewidth - 4\tabcolsep) * \real{0.3333}}
  >{\raggedright\arraybackslash}p{(\linewidth - 4\tabcolsep) * \real{0.3333}}
  >{\raggedright\arraybackslash}p{(\linewidth - 4\tabcolsep) * \real{0.3333}}@{}}
\toprule\noalign{}
\begin{minipage}[b]{\linewidth}\raggedright
Transformer 개념
\end{minipage} & \begin{minipage}[b]{\linewidth}\raggedright
무선통신/신호처리 대응
\end{minipage} & \begin{minipage}[b]{\linewidth}\raggedright
설명
\end{minipage} \\
\midrule\noalign{}
\endhead
\bottomrule\noalign{}
\endlastfoot
쿼리·키 내적 \(q\cdot k\) & \textbf{정합 필터(matched filter)/상관기} &
수신 신호와 템플릿의 상관값. 어텐션 점수는 곧 상관값이다. \\
코사인 유사도 & \textbf{정규화 상관계수} & 크기를 제거하고 방향(파형
모양)만 비교 --- 신호 검출의 정규화 상관과 동일. \\
소프트맥스 가중합 & \textbf{최대비결합(MRC)류 소프트 결합} &
SNR/유사도가 높은 경로에 큰 가중치를 주는 가중 결합. \\
고정 가중치 vs 어텐션 & \textbf{고정 빔포밍 vs 적응 빔포밍(MVDR/LCMV)} &
고정 빔포머는 입력과 무관한 고정 가중치, 적응 빔포머는 \textbf{수신
데이터에 맞춰 가중치를 갱신} --- 어텐션이 정확히 후자다. \\
멀티헤드 & \textbf{MIMO 공간 다중화/다중 빔} & 서로 다른 헤드 = 서로
다른 공간 채널/빔을 동시에 형성. \\
위치 인코딩(사인파) & \textbf{OFDM 부반송파/파일럿 인덱싱} & 서로 다른
주파수의 사인파로 ``위치''를 직교적으로 부호화 --- OFDM이 부반송파
주파수로 자원을 인덱싱하는 것과 형식이 같다. \\
KV 캐시 & \textbf{채널 상태/지연 프로파일 메모리} & 과거 상태를 저장해
긴 시간축 의존성을 유지. \\
\(O(n^2)\) 복잡도 & \textbf{전 쌍 상관행렬 계산 비용} & 모든 쿼리-키
쌍의 상관 = 공분산 행렬 규모의 연산. \\
\end{longtable}
}

\begin{tcolorbox}[enhanced jigsaw, arc=.35mm, bottomrule=.15mm, bottomtitle=1mm, breakable, colback=white, colbacktitle=quarto-callout-tip-color!10!white, colframe=quarto-callout-tip-color-frame, coltitle=black, left=2mm, leftrule=.75mm, opacityback=0, opacitybacktitle=0.6, rightrule=.15mm, title=\textcolor{quarto-callout-tip-color}{\faLightbulb}\hspace{0.5em}{한 줄 요약 비유}, titlerule=0mm, toprule=.15mm, toptitle=1mm]

\textbf{어텐션 ≈ 데이터에 적응하는 빔포밍.} 고정 빔포머(보통의 가중치)는
어디서 신호가 오든 같은 가중치를 쓰지만, 적응 빔포머(어텐션)는 수신
데이터(쿼리·키)를 보고 가중치를 재설계해 원하는 방향(유사한 토큰)으로
빔을 형성한다.

\end{tcolorbox}

\begin{center}\rule{0.5\linewidth}{0.5pt}\end{center}

\section{핵심 요약}\label{sec-summary}

{\def\LTcaptype{none} % do not increment counter
\begin{longtable}[]{@{}
  >{\raggedright\arraybackslash}p{(\linewidth - 2\tabcolsep) * \real{0.5000}}
  >{\raggedright\arraybackslash}p{(\linewidth - 2\tabcolsep) * \real{0.5000}}@{}}
\toprule\noalign{}
\begin{minipage}[b]{\linewidth}\raggedright
주제
\end{minipage} & \begin{minipage}[b]{\linewidth}\raggedright
반드시 기억할 것
\end{minipage} \\
\midrule\noalign{}
\endhead
\bottomrule\noalign{}
\endlastfoot
세 기둥 & 토큰 · 어텐션 · 위치 인코딩 (어텐션이 가장 중요) \\
어텐션의 본질 & 정보가 아니라 \textbf{유사도에 학습된 가중치}를 건다;
가중치가 \textbf{입력의 함수} \\
Q/K/V & 쿼리=물어보는 것, 키=비교 상대, 값=가져올 정보; 학습 대상은
\(W_q,W_k,W_v\) \\
유사도 & 내적 = 코사인 유사도(방향 비교) → 소프트맥스로 차이 증폭 \\
스케일링 & \(\sqrt{d_k}\) 로 나눠 소프트맥스 과포화/기울기 소실 방지 \\
셀프 어텐션 & Q,K,V 동일 출처; 문맥 반영해 표현 정련 \\
마스크드 어텐션 & 디코더에서 미래 토큰 차단(인과성); 잘못된 기울기 전파
방지 \\
멀티헤드 & 관계의 여러 측면; CNN의 다중 필터와 같은 동기 \\
순열 등변성 & 어텐션은 순서를 모른다 → 위치 인코딩 필요 \\
응용 & GPT(=Generative Pretrained Transformer), ViT \\
한계 & 윈도 밖 문맥 손실, \(O(n^2)\); 완화책 KV 캐시, 대안 Mamba \\
\end{longtable}
}

\begin{center}\rule{0.5\linewidth}{0.5pt}\end{center}

\section{부록: NumPy 구현}\label{sec-appendix}

아래 코드는 어텐션, 인과 마스크, 멀티헤드, 사인파 위치 인코딩의 동작을
최소 구현으로 확인한다.

\begin{Shaded}
\begin{Highlighting}[]
\ImportTok{import}\NormalTok{ numpy }\ImportTok{as}\NormalTok{ np}

\KeywordTok{def}\NormalTok{ softmax(x, axis}\OperatorTok{={-}}\DecValTok{1}\NormalTok{):}
\NormalTok{    x }\OperatorTok{=}\NormalTok{ x }\OperatorTok{{-}}\NormalTok{ np.}\BuiltInTok{max}\NormalTok{(x, axis}\OperatorTok{=}\NormalTok{axis, keepdims}\OperatorTok{=}\VariableTok{True}\NormalTok{)}
\NormalTok{    e }\OperatorTok{=}\NormalTok{ np.exp(x)}
    \ControlFlowTok{return}\NormalTok{ e }\OperatorTok{/}\NormalTok{ np.}\BuiltInTok{sum}\NormalTok{(e, axis}\OperatorTok{=}\NormalTok{axis, keepdims}\OperatorTok{=}\VariableTok{True}\NormalTok{)}

\KeywordTok{def}\NormalTok{ scaled\_dot\_product\_attention(Q, K, V, mask}\OperatorTok{=}\VariableTok{None}\NormalTok{):}
    \CommentTok{"""Attention(Q,K,V) = softmax(QK\^{}T / sqrt(d\_k)) V"""}
\NormalTok{    d\_k }\OperatorTok{=}\NormalTok{ Q.shape[}\OperatorTok{{-}}\DecValTok{1}\NormalTok{]}
\NormalTok{    scores }\OperatorTok{=}\NormalTok{ Q }\OperatorTok{@}\NormalTok{ np.swapaxes(K, }\OperatorTok{{-}}\DecValTok{1}\NormalTok{, }\OperatorTok{{-}}\DecValTok{2}\NormalTok{) }\OperatorTok{/}\NormalTok{ np.sqrt(d\_k)   }\CommentTok{\# 쿼리{-}키 상관 (정합필터)}
    \ControlFlowTok{if}\NormalTok{ mask }\KeywordTok{is} \KeywordTok{not} \VariableTok{None}\NormalTok{:}
\NormalTok{        scores }\OperatorTok{=}\NormalTok{ np.where(mask, scores, }\OperatorTok{{-}}\FloatTok{1e9}\NormalTok{)            }\CommentTok{\# 마스킹 위치는 {-}inf}
\NormalTok{    A }\OperatorTok{=}\NormalTok{ softmax(scores, axis}\OperatorTok{={-}}\DecValTok{1}\NormalTok{)                         }\CommentTok{\# 어텐션 가중치 (합=1)}
    \ControlFlowTok{return}\NormalTok{ A }\OperatorTok{@}\NormalTok{ V, A}
\end{Highlighting}
\end{Shaded}

\begin{Shaded}
\begin{Highlighting}[]
\CommentTok{\# 1) 셀프 어텐션}
\NormalTok{rng }\OperatorTok{=}\NormalTok{ np.random.default\_rng(}\DecValTok{0}\NormalTok{)}
\NormalTok{n, d }\OperatorTok{=} \DecValTok{4}\NormalTok{, }\DecValTok{8}
\NormalTok{X }\OperatorTok{=}\NormalTok{ rng.standard\_normal((n, d))}
\NormalTok{Wq, Wk, Wv }\OperatorTok{=}\NormalTok{ (rng.standard\_normal((d, d)) }\ControlFlowTok{for}\NormalTok{ \_ }\KeywordTok{in} \BuiltInTok{range}\NormalTok{(}\DecValTok{3}\NormalTok{))}
\NormalTok{Q, K, V }\OperatorTok{=}\NormalTok{ X }\OperatorTok{@}\NormalTok{ Wq, X }\OperatorTok{@}\NormalTok{ Wk, X }\OperatorTok{@}\NormalTok{ Wv}
\NormalTok{out, A }\OperatorTok{=}\NormalTok{ scaled\_dot\_product\_attention(Q, K, V)}
\ControlFlowTok{assert}\NormalTok{ np.allclose(A.}\BuiltInTok{sum}\NormalTok{(axis}\OperatorTok{=}\DecValTok{1}\NormalTok{), }\FloatTok{1.0}\NormalTok{)        }\CommentTok{\# 각 행의 가중치 합 = 1}

\CommentTok{\# 2) 인과(causal) 마스크: 하삼각만 허용 → 미래 차단}
\NormalTok{causal }\OperatorTok{=}\NormalTok{ np.tril(np.ones((n, n), dtype}\OperatorTok{=}\BuiltInTok{bool}\NormalTok{))}
\NormalTok{out\_c, A\_c }\OperatorTok{=}\NormalTok{ scaled\_dot\_product\_attention(Q, K, V, mask}\OperatorTok{=}\NormalTok{causal)}
\ControlFlowTok{assert}\NormalTok{ np.allclose(A\_c[np.triu\_indices(n, }\DecValTok{1}\NormalTok{)], }\DecValTok{0}\NormalTok{, atol}\OperatorTok{=}\FloatTok{1e{-}6}\NormalTok{)  }\CommentTok{\# 미래 가중치 \textasciitilde{} 0}
\end{Highlighting}
\end{Shaded}

\begin{Shaded}
\begin{Highlighting}[]
\CommentTok{\# 3) 멀티헤드 어텐션}
\KeywordTok{class}\NormalTok{ MultiHeadAttention:}
    \KeywordTok{def} \FunctionTok{\_\_init\_\_}\NormalTok{(}\VariableTok{self}\NormalTok{, d\_model, h, seed}\OperatorTok{=}\DecValTok{0}\NormalTok{):}
\NormalTok{        rng }\OperatorTok{=}\NormalTok{ np.random.default\_rng(seed)}
        \VariableTok{self}\NormalTok{.h, }\VariableTok{self}\NormalTok{.dk }\OperatorTok{=}\NormalTok{ h, d\_model }\OperatorTok{//}\NormalTok{ h}
        \VariableTok{self}\NormalTok{.Wq, }\VariableTok{self}\NormalTok{.Wk, }\VariableTok{self}\NormalTok{.Wv, }\VariableTok{self}\NormalTok{.Wo }\OperatorTok{=}\NormalTok{ (}
\NormalTok{            rng.standard\_normal((d\_model, d\_model)) }\ControlFlowTok{for}\NormalTok{ \_ }\KeywordTok{in} \BuiltInTok{range}\NormalTok{(}\DecValTok{4}\NormalTok{))}
    \KeywordTok{def}\NormalTok{ \_split(}\VariableTok{self}\NormalTok{, x):                 }\CommentTok{\# (n, d) {-}\textgreater{} (h, n, dk)}
\NormalTok{        n, \_ }\OperatorTok{=}\NormalTok{ x.shape}
        \ControlFlowTok{return}\NormalTok{ x.reshape(n, }\VariableTok{self}\NormalTok{.h, }\VariableTok{self}\NormalTok{.dk).transpose(}\DecValTok{1}\NormalTok{, }\DecValTok{0}\NormalTok{, }\DecValTok{2}\NormalTok{)}
    \KeywordTok{def} \FunctionTok{\_\_call\_\_}\NormalTok{(}\VariableTok{self}\NormalTok{, X, mask}\OperatorTok{=}\VariableTok{None}\NormalTok{):}
\NormalTok{        Q, K, V }\OperatorTok{=} \VariableTok{self}\NormalTok{.\_split(X }\OperatorTok{@} \VariableTok{self}\NormalTok{.Wq), }\VariableTok{self}\NormalTok{.\_split(X }\OperatorTok{@} \VariableTok{self}\NormalTok{.Wk), }\VariableTok{self}\NormalTok{.\_split(X }\OperatorTok{@} \VariableTok{self}\NormalTok{.Wv)}
\NormalTok{        out, \_ }\OperatorTok{=}\NormalTok{ scaled\_dot\_product\_attention(Q, K, V, mask)   }\CommentTok{\# (h, n, dk)}
\NormalTok{        out }\OperatorTok{=}\NormalTok{ out.transpose(}\DecValTok{1}\NormalTok{, }\DecValTok{0}\NormalTok{, }\DecValTok{2}\NormalTok{).reshape(X.shape[}\DecValTok{0}\NormalTok{], }\OperatorTok{{-}}\DecValTok{1}\NormalTok{)   }\CommentTok{\# concat heads}
        \ControlFlowTok{return}\NormalTok{ out }\OperatorTok{@} \VariableTok{self}\NormalTok{.Wo}

\NormalTok{y }\OperatorTok{=}\NormalTok{ MultiHeadAttention(d\_model}\OperatorTok{=}\DecValTok{16}\NormalTok{, h}\OperatorTok{=}\DecValTok{4}\NormalTok{)(rng.standard\_normal((}\DecValTok{5}\NormalTok{, }\DecValTok{16}\NormalTok{)))}
\ControlFlowTok{assert}\NormalTok{ y.shape }\OperatorTok{==}\NormalTok{ (}\DecValTok{5}\NormalTok{, }\DecValTok{16}\NormalTok{)}
\end{Highlighting}
\end{Shaded}

\begin{Shaded}
\begin{Highlighting}[]
\CommentTok{\# 4) 사인파 위치 인코딩}
\KeywordTok{def}\NormalTok{ positional\_encoding(seq\_len, d\_model):}
\NormalTok{    pos }\OperatorTok{=}\NormalTok{ np.arange(seq\_len)[:, }\VariableTok{None}\NormalTok{]}
\NormalTok{    i }\OperatorTok{=}\NormalTok{ np.arange(d\_model)[}\VariableTok{None}\NormalTok{, :]}
\NormalTok{    angle }\OperatorTok{=}\NormalTok{ pos }\OperatorTok{/}\NormalTok{ np.power(}\DecValTok{10000}\NormalTok{, (}\DecValTok{2} \OperatorTok{*}\NormalTok{ (i }\OperatorTok{//} \DecValTok{2}\NormalTok{)) }\OperatorTok{/}\NormalTok{ d\_model)}
    \ControlFlowTok{return}\NormalTok{ np.where(i }\OperatorTok{\%} \DecValTok{2} \OperatorTok{==} \DecValTok{0}\NormalTok{, np.sin(angle), np.cos(angle))}

\NormalTok{PE }\OperatorTok{=}\NormalTok{ positional\_encoding(seq\_len}\OperatorTok{=}\DecValTok{10}\NormalTok{, d\_model}\OperatorTok{=}\DecValTok{16}\NormalTok{)}
\ControlFlowTok{assert}\NormalTok{ PE.shape }\OperatorTok{==}\NormalTok{ (}\DecValTok{10}\NormalTok{, }\DecValTok{16}\NormalTok{) }\KeywordTok{and}\NormalTok{ PE.}\BuiltInTok{min}\NormalTok{() }\OperatorTok{\textgreater{}=} \OperatorTok{{-}}\DecValTok{1} \KeywordTok{and}\NormalTok{ PE.}\BuiltInTok{max}\NormalTok{() }\OperatorTok{\textless{}=} \DecValTok{1}
\BuiltInTok{print}\NormalTok{(}\StringTok{"모든 검증 통과"}\NormalTok{)}
\end{Highlighting}
\end{Shaded}

\begin{center}\rule{0.5\linewidth}{0.5pt}\end{center}

\section{더 읽을거리}\label{sec-reading}

\begin{itemize}
\tightlist
\item
  d2l.ai --- \emph{Attention Mechanisms and Transformers} 장: 어텐션
  풀링, 스코어 함수(\(\sqrt{d_k}\) 스케일링), 멀티헤드, 셀프 어텐션과
  위치 인코딩, Transformer 아키텍처를 코드와 함께 다룬다.
  \url{https://d2l.ai/}
\item
  \emph{Mathematics for Machine Learning} (G. Thomas) ---
  내적·사영·각거리(코사인 유사도)의 기초가 되는 선형대수, 그리고
  소프트맥스/확률의 통계적 배경.
  \url{https://gwthomas.github.io/docs/math4ml.pdf}
\end{itemize}




\end{document}
